\documentclass[11pt,a4paper]{article}
\usepackage[utf8]{inputenc}
\usepackage[spanish]{babel}
\usepackage{amsmath, amssymb}
\usepackage{graphicx}
\usepackage[margin=2.5cm]{geometry}
\usepackage{fancyhdr}
\usepackage{tcolorbox}
\usepackage{enumitem}

% Configuración de colores
\definecolor{UCBblue}{RGB}{0, 51, 153}
\definecolor{UCBgray}{RGB}{100, 100, 100}

% Configuración de encabezado y pie de página
\pagestyle{fancy}
\fancyhf{}
\fancyhead[L]{\textcolor{UCBblue}{\textbf{Ingeniería Mecatrónica - UCB Tarija}}}
\fancyhead[R]{\textcolor{UCBgray}{Taller de Nivelación - Sesión 1}}
\fancyfoot[C]{\thepage}
\renewcommand{\headrulewidth}{0.4pt}
\renewcommand{\headrule}{\hbox to\headwidth{\color{UCBblue}\leaders\hrule height \headrulewidth\hfill}}

\begin{document}

\begin{center}
    {\Large \textbf{\textcolor{UCBblue}{Misión de Ingeniería 1: Modelado de Redes de Potencia}}}\\
    \vspace{0.3cm}
    {\large \textbf{Fase CDIO: Concebir (Just-In-Time Physics)}}\\
    \vspace{0.5cm}
    \textbf{Estudiante:} \rule{8cm}{0.4pt} \quad \textbf{Fecha:} \rule{3cm}{0.4pt}
\end{center}

\vspace{0.5cm}

\section*{1. El Concepto Físico: Entendiendo la Electricidad}
Antes de escribir código, debemos entender la física del problema. Imagina que el circuito eléctrico es un sistema de tuberías de agua en una planta industrial:
\begin{itemize}
    \item \textbf{Voltaje (V):} Es la bomba de agua. Aplica presión para empujar el fluido.
    \item \textbf{Corriente (I):} Es el caudal de agua fluyendo por las tuberías (Amperios).
    \item \textbf{Resistencia (R):} Son estrechamientos en la tubería que dificultan el paso del agua (Ohmios).
    \item \textbf{Malla:} Es cualquier circuito cerrado (como una pista de carreras) que el agua puede recorrer.
\end{itemize}

\section*{2. La Receta Algorítmica: El Método de Mallas (LVK)}
Para programar simuladores, los ingenieros no resuelven ecuaciones al azar; siguen un algoritmo estricto para construir una \textbf{Matriz de Resistencias ($\mathbf{R}$)}. Sigue estos 3 pasos infalibles:

\begin{tcolorbox}[colback=blue!5!white,colframe=UCBblue,title=\textbf{El Algoritmo Visual para Matrices}]
\textbf{Paso 1: Asignar Sentidos} \\
Dibuja una flecha circular en sentido horario ($\circlearrowright$) dentro de cada malla. Nómbralas $I_1, I_2, I_3$.

\vspace{0.2cm}
\textbf{Paso 2: La Diagonal Principal (Resistencias Propias)} \\
Para encontrar el valor de $R_{11}$ (Malla 1), simplemente suma \textbf{todas} las resistencias que toca la flecha $I_1$. Repite para $R_{22}$ y $R_{33}$. \textit{Regla de oro: Estos valores siempre son positivos.}

\vspace{0.2cm}
\textbf{Paso 3: Fuera de la Diagonal (Resistencias Mutuas)} \\
Mira la "frontera" entre dos mallas. La resistencia que está justo al medio es la compartida. Si buscas $R_{12}$, anota la resistencia entre la Malla 1 y la Malla 2, pero \textbf{cámbiale el signo a negativo}. \textit{Regla de oro: En redes pasivas, $R_{12}$ es exactamente igual a $R_{21}$ (Simetría).}
\end{tcolorbox}

\section*{3. Tratamiento de las Fuentes de Poder (Vector $\mathbf{V}$)}
Para construir el lado derecho de nuestra ecuación ($\mathbf{R}\mathbf{I} = \mathbf{V}$), sigue el giro de tu flecha de corriente:
\begin{itemize}
    \item Si la flecha choca primero con la línea corta (negativo) del símbolo de la batería, el voltaje "sube", por lo que se anota \textbf{positivo} en el vector.
    \item Si choca primero con la línea larga (positivo), el voltaje cae, se anota \textbf{negativo}.
\end{itemize}

\newpage

\section*{4. Misión Práctica: Desarrollo Manuscrito}

\begin{figure}[h!]
    \centering
    % Placeholder para que el docente pegue la imagen exportada de LTSpice
    \framebox[12cm][c]{\rule{0pt}{5cm} \textit{Espacio reservado para el plano esquemático del actuador (Docente: Insertar imagen LTSpice aquí)}}
    \caption{Plano de distribución de potencia de 3 mallas.}
\end{figure}

Dado el circuito superior con los siguientes valores de componentes:
$$R_1=10\Omega, \quad R_2=22\Omega, \quad R_3=15\Omega, \quad R_4=33\Omega, \quad R_5=47\Omega, \quad R_6=12\Omega$$
$$V_{s1} = 24V \text{ (Malla 1)}, \quad V_{s2} = 0V \text{ (Malla 2)}, \quad V_{s3} = -12V \text{ (Malla 3)}$$

\textbf{A.} Utilizando "La Receta Algorítmica", llena los espacios para estructurar el sistema canónico $\mathbf{R}\mathbf{I} = \mathbf{V}$:

\vspace{0.5cm}

$$
\begin{bmatrix}
 (R_1 + \dots \dots + R_3) & -R_2 & -R_3 \\
 \dots \dots \dots \dots \dots & (\dots \dots + R_4 + \dots \dots) & -R_5 \\
 \dots \dots \dots \dots \dots & \dots \dots \dots \dots \dots & (R_3 + \dots \dots + R_6)
\end{bmatrix}
\begin{bmatrix}
 I_1 \\
 I_2 \\
 I_3
\end{bmatrix}
=
\begin{bmatrix}
 24.0 \\
 \dots \dots \\
 \dots \dots
\end{bmatrix}
$$

\vspace{0.5cm}

\textbf{B.} Reemplaza los valores numéricos y calcula a mano el determinante de la matriz $\mathbf{R}$ para diagnosticar si existe un cortocircuito ($\det(\mathbf{R}) = 0$). Usa la Regla de Sarrus o cofactores.

\vspace{0.3cm}
\begin{tcolorbox}[colback=white,colframe=UCBgray,height=6cm,valign=top]
\textit{Espacio para desarrollo del determinante (Sarrus):}
\end{tcolorbox}

\section*{5. Operar: Hacia el Gemelo Digital}
Una vez que hayas certificado tu modelo matemático manuscrito:
\begin{enumerate}[label=\textbf{\arabic*.}]
    \item Abre tu entorno de programación (Jupyter).
    \item Ejecuta el \texttt{Notebook\_01\_Mallas.ipynb}.
    \item Transcribe los valores de tu matriz $\mathbf{R}$ al código Python.
    \item Verifica que la celda de \texttt{assert} apruebe tu matriz. ¡Si es simétrica, has superado la fase de Diseño!
\end{enumerate}

\end{document}
